% Reviewed translation: PDF pages 346--350(part) / printed pages 332--336(part).
\MathBookSourcePage{346}
现在可以定义对偶线性化 Navier--Stokes 系统:
对 $\boldsymbol g\in L^2(\Omega)^2$, 求
$z=(\operatorname{curl}\chi,\mu)\in V$, 使得
\begin{equation}\label{eq:navier-centered-stream-dual-linearized}
\widetilde a((I+TD^*)z,v)
=(\boldsymbol g,\operatorname{curl}\phi)
\qquad
\forall v=(\operatorname{curl}\phi,\theta)\in V;
\end{equation}
即
\[
z=(I+TD^*)^{-1}T\boldsymbol g.
\]

\setcounter{statement}{0}
\begin{Lemma}\label{lem:navier-centered-stream-dual-regularity}
设 $\Omega$ 是 $\mathbb R^2$ 中有界、Lipschitz 连续的区域,
并假设 Stokes 问题具有正则性
\eqref{eq:navier-centered-gp-regularity}. 则 Problem~
\eqref{eq:navier-centered-stream-dual-linearized} 的解
$z=(\operatorname{curl}\chi,\mu)$ 属于
$H^2(\Omega)^2\times H^1(\Omega)$, 且
\begin{equation}\label{eq:navier-centered-stream-dual-regularity}
\|\chi\|_{3,\Omega}+\|\mu\|_{1,\Omega}
\leq C\{1+\lambda(\|\psi(\lambda)\|_{3,\Omega}
+\|\omega(\lambda)\|_{1,\Omega})\}
\|\boldsymbol g\|_{0,\Omega}.
\end{equation}
\end{Lemma}
\begin{Proof}
首先由 \eqref{eq:navier-centered-stream-dual-linearized} 得
\begin{equation}\label{eq:navier-centered-stream-dual-bound}
\|z\|_X\leq C_1\|\boldsymbol g\|_{0,\Omega}.
\end{equation}
接着, Lemma~\ref{lem:incompressible-weak-kernel-coincidence} 告诉我们,
所有 $v=(\operatorname{curl}\phi,\theta)\in V$ 都满足
\[
\operatorname{curl}\phi\in H_0^1(\Omega)^2,
\qquad
\theta=-\Delta\phi,
\qquad
\|\theta\|_{0,\Omega}=|v|\simeq\|v\|_X.
\]
因此
\[
\begin{aligned}
\langle D^*z,v\rangle
&=\lambda(\operatorname{curl}\chi,
\omega(\lambda)\operatorname{grad}\phi
-\Delta\phi\operatorname{grad}\psi(\lambda))\\
&=\lambda(\operatorname{curl}
(\operatorname{curl}\chi\cdot\operatorname{grad}\psi(\lambda))
-\omega(\lambda)\operatorname{grad}\chi,
\operatorname{curl}\phi).
\end{aligned}
\]
于是 $D^*z$ 可写成
\[
\langle D^*z,v\rangle=(\boldsymbol l,\operatorname{curl}\phi)
\qquad\forall v=(\operatorname{curl}\phi,\theta)\in V,
\]
其中由 $u(\lambda)$ 的正则性和 Sobolev Embedding
Theorem~\ref{thm:sobolev-embedding}, 有 $\boldsymbol l\in L^2(\Omega)^2$, 且
\[
\begin{aligned}
\|\boldsymbol l\|_{0,\Omega}
&\leq C_2\lambda(\|\omega(\lambda)\|_{1,\Omega}
+\|\psi(\lambda)\|_{3,\Omega})\|\chi\|_{2,\Omega}\\
&\leq C_3\lambda(\|\omega(\lambda)\|_{1,\Omega}
+\|\psi(\lambda)\|_{3,\Omega})
\|\boldsymbol g\|_{0,\Omega},
\end{aligned}
\]
因为 $z\in V$ 且满足 \eqref{eq:navier-centered-stream-dual-bound}.
最后一个不等式、正则性假设
\eqref{eq:navier-centered-gp-regularity}, 以及 $z$ 还可表示为
$z=T(\boldsymbol g-D^*z)$ 这一事实给出
\eqref{eq:navier-centered-stream-dual-regularity}.
\end{Proof}

\setcounter{statement}{4}
\begin{Theorem}\label{thm:navier-centered-stream-h1-error}
设 $\Omega$ 和 $\mathcal T_h$ 如 Theorem~
\ref{thm:navier-centered-stream-branch}, 并假设 Navier--Stokes
Problem~\eqref{eq:navier-centered-stream-dual-bound} 有一个非奇异解分支,
使映射 $\lambda\mapsto\psi(\lambda)$ 在 $l\geq2$ 时从 $\Lambda$ 到
$H^{l+1}(\Omega)$ 连续, 在 $l=1$ 时从 $\Lambda$ 到 $H^3(\Omega)$ 连续.
则对所有 $\lambda\in\Lambda$ 有估计
\begin{equation}\label{eq:navier-centered-stream-h1-error}
|\psi(\lambda)-\psi_h(\lambda)|_{1,\Omega}
\leq
\begin{cases}
C_1h^l,&l\geq2,\\
C_2(\varepsilon)h^{1-\varepsilon},&l=1,
\quad\forall\varepsilon>0,
\end{cases}
\end{equation}
其中常数与 $h$ 和 $\lambda$ 无关.
\end{Theorem}

\MathBookSourcePage{347}
\begin{Proof}
令 $\boldsymbol g$ 是 $L^2(\Omega)^2$ 中的任意函数, 并令
$z=(\operatorname{curl}\chi,\nu)\in V$ 是线性化问题
\eqref{eq:navier-centered-stream-dual-linearized} 的相应解:
\[
(\boldsymbol g,\operatorname{curl}\phi)
=\widetilde a(z,v)+(\operatorname{curl}\chi,D\cdot v)
+\widetilde b(v,\nu)
\qquad
\forall v=(\operatorname{curl}\phi,\theta)\in X.
\]
特别地(暂时略去参数 $\lambda$),
\begin{equation}\label{eq:navier-centered-stream-dual-error-identity}
\begin{aligned}
(\boldsymbol g,\operatorname{curl}(\psi-\psi_h))
={}&\widetilde a(z,u-u_h)
+(\operatorname{curl}\chi,D\cdot(u-u_h))\\
&+\widetilde b(u-u_h,\nu).
\end{aligned}
\end{equation}
但由 \eqref{eq:navier-centered-stream-continuous-vorticity} 和
\eqref{eq:navier-centered-stream-discrete-vorticity}, $u-u_h$ 满足
\begin{equation}\label{eq:navier-centered-stream-galerkin-error}
\widetilde a(u-u_h,z_h)+\widetilde b(z_h,\omega)
=-(G(\lambda,u)-G(\lambda,u_h),\operatorname{curl}\chi_h)
\end{equation}
对每个 $z_h=(\operatorname{curl}\chi_h,\nu_h)\in V_h$ 都成立, 其中
\[
V_h=\{v_h=(\operatorname{curl}\phi_h,\theta_h);
\ \phi_h\in\Phi_h,\ \theta_h\in\Theta_h,
\ \widetilde b(v_h,\mu_h)=0
\quad\forall\mu_h\in\Theta_h\}.
\]
因此, 合并 \eqref{eq:navier-centered-stream-dual-error-identity} 和
\eqref{eq:navier-centered-stream-galerkin-error}, 并利用
$u,z\in V$ 以及 $u_h,z_h\in V_h$, 得
\[
\begin{aligned}
(\boldsymbol g,\operatorname{curl}(\psi-\psi_h))
={}&\widetilde a(z-z_h,u-u_h)
+(\operatorname{curl}(\chi-\chi_h),D\cdot(u-u_h))\\
&+\widetilde b(u-u_h,\nu-\theta_h)
+\widetilde b(z-z_h,\omega-\mu_h)\\
&+(\operatorname{curl}\chi_h,
D\cdot(u-u_h)-(G(\lambda,u)-G(\lambda,u_h)))
\end{aligned}
\]
对任意 $z_h\in V_h$ 以及任意 $\mu_h,\theta_h\in\Theta_h$ 均成立.

取 $\theta_h=P_h\nu$ 和 $\mu_h=P_h\omega$;
公式 \eqref{eq:appendix-dirichlet-ritz-projection} 对所有
$\phi_h,\delta_h\in\Phi_h$ 给出
\[
\widetilde b(u-u_h,\nu-\theta_h)
=(\operatorname{curl}(\psi-\phi_h),
\operatorname{curl}(\nu-P_h\nu))
-(\omega-\omega_h,\nu-P_h\nu),
\]
\[
\widetilde b(z-z_h,\omega-\mu_h)
=(\operatorname{curl}(\chi-\delta_h),
\operatorname{curl}(\omega-P_h\omega))
-(\nu-\nu_h,\omega-P_h\omega).
\]
另一方面, Taylor 公式 \eqref{eq:navier-branch-taylor-remainder} 在这里给出
\[
\begin{aligned}
G(\lambda,u)-G(\lambda,u_h)
-D_uG(\lambda,u)\cdot(u-u_h)
&=-\frac12D_{uu}^2G\cdot(u-u_h)^2\\
&=-\lambda(\omega-\omega_h)
\operatorname{grad}(\psi-\psi_h).
\end{aligned}
\]
因此, 对所有 $z_h\in V_h$ 以及 $\phi_h,\delta_h\in\Phi_h$, 有
\begin{equation}\label{eq:navier-centered-stream-expanded-duality}
\begin{aligned}
(\boldsymbol g,\operatorname{curl}(\psi-\psi_h))
={}&\widetilde a(z-z_h,u-u_h)\\
&+\lambda(\operatorname{curl}(\chi-\chi_h),
\omega\operatorname{grad}(\psi-\psi_h))\\
&+\lambda(\operatorname{curl}(\chi-\chi_h),
(\omega-\omega_h)\operatorname{grad}\psi)\\
&+\lambda(\operatorname{curl}\chi_h,
(\omega-\omega_h)\operatorname{grad}(\psi-\psi_h))\\
&+(\operatorname{curl}(\psi-\phi_h),
\operatorname{curl}(\nu-P_h\nu))
-(\omega-\omega_h,\nu-P_h\nu)\\
&+(\operatorname{curl}(\chi-\delta_h),
\operatorname{curl}(\omega-P_h\omega))
-(\nu-\nu_h,\omega-P_h\omega).
\end{aligned}
\end{equation}

\MathBookSourcePage{348}
由此立即推出
\begin{equation}\label{eq:navier-centered-stream-duality-bound}
\begin{aligned}
|(&\boldsymbol g,\operatorname{curl}(\psi-\psi_h))|
\leq\|u-u_h\|_X
\{\,\|z-z_h\|_X
+\lambda|\psi|_{1,4,\Omega}|\chi-\chi_h|_{1,4,\Omega}\\
&+\lambda|\chi_h|_{1,4,\Omega}|\psi-\psi_h|_{1,4,\Omega}
+\|\nu-P_h\nu\|_{0,\Omega}\}\\
&+\lambda\|u\|_X|\psi-\psi_h|_{1,4,\Omega}
|\chi-\chi_h|_{1,4,\Omega}\\
&+|\nu-P_h\nu|_{1,\Omega}|\psi-\phi_h|_{1,\Omega}\\
&+|\chi-\delta_h|_{1,\Omega}|\omega-P_h\omega|_{1,\Omega}\\
&+\|z-z_h\|_X\|\omega-P_h\omega\|_{0,\Omega},
\qquad
\forall z_h\in V_h,
\quad\forall\phi_h,\delta_h\in\Phi_h.
\end{aligned}
\end{equation}
最后回忆(参见 Lemma~
\ref{lem:incompressible-further-improved-kernel-approximation} 和
Remark~\ref{rem:incompressible-further-near-smooth-error}):
\[
\inf_{z_h\in V_h}\|z-z_h\|_X
\leq
\begin{cases}
C_1h\|\chi\|_{3,\Omega},&l\geq2,\\
C_2(\varepsilon)h^{1/2-\varepsilon}\|\chi\|_{3,\Omega},
&l=1,\quad\forall\varepsilon>0;
\end{cases}
\]
另一方面,
\[
\inf_{\delta_h\in\Phi_h}|\chi-\delta_h|_{1,\Omega}
\leq C_3h^\alpha|\psi|_{\alpha+1,\Omega},
\qquad
\alpha=\min(l,2);
\]
类似地,
\[
\inf_{\phi_h\in\Phi_h}|\psi-\phi_h|_{1,\Omega}
\leq C_4h^l|\psi|_{l+1,\Omega},
\]
以及
\[
\|\omega-P_h\omega\|_{0,\Omega}
+h|\omega-P_h\omega|_{1,\Omega}
\leq C_5h^\beta|\omega|_{\beta,\Omega},
\]
其中 $l=1$ 时 $\beta=1$, $l\geq2$ 时 $\beta=l-1$.
把这些界代入 \eqref{eq:navier-centered-stream-duality-bound},
并应用 Theorem~\ref{thm:navier-centered-stream-branch} 和
Lemma~\ref{lem:navier-centered-stream-dual-regularity}, 即容易得到
\eqref{eq:navier-centered-stream-h1-error}.
\end{Proof}

\subsection{关于``流函数--速度梯度张量''格式的评注}
\label{subsec:navier-centered-stream-gradient-tensor}

对 Section~\ref{subsec:navier-centered-stream-vorticity} 的方法作少量修改,
就可以把它应用于 Navier--Stokes 方程的``流函数--速度梯度张量''方法,
至少在 $\Omega$ 是平面凸多边形时如此. 回到
Subsection~\ref{subsec:incompressible-hhj-formulation}, 回忆双线性形式
\[
a_h(\boldsymbol\sigma,\boldsymbol\tau)
=\int_\Omega\sigma_{ij}\tau_{ij}\,dx,
\]
\[
b_h(\boldsymbol\tau,\phi)
=-\sum_{\kappa\in\mathcal T_h}\int_\kappa
\tau_{ij}\frac{\partial^2\phi}{\partial x_i\partial x_j}\,dx
+\int_{\Gamma_h}M_n(\boldsymbol\tau)
S\left(\frac{\partial\phi}{\partial n}\right)ds,
\]
以及空间
\[
\begin{aligned}
\widetilde\Sigma
&=\{\boldsymbol\tau=(\tau_{ij})\in L^2(\Omega)^4;
\ \tau_{12}=\tau_{21},
\ \boldsymbol\tau|_\kappa\in W^{1,r}(\kappa)^4
\quad\forall\kappa\in\mathcal T_h,\\
&\hspace{4em}
M_n(\boldsymbol\tau)\text{ is continuous on each segment of }\Gamma_h\},
\end{aligned}
\]
\[
\widetilde\Psi
=\{\phi\in\Phi_s;
\ \phi|_\kappa\in H^2(\kappa)
\quad\forall\kappa\in\mathcal T_h\},
\]

\MathBookSourcePage{349}
其中固定 $s\geq4$ 且 $1/r+1/s=1$.
对 $\boldsymbol f\in L^r(\Omega)^2$, 已知 Stokes 算子可以表示为
\[
\left\{
\begin{aligned}
T\boldsymbol f&=u=(\boldsymbol\sigma,\psi)
\in\widetilde\Sigma\times\widetilde\Psi,\\
b_h(\boldsymbol\sigma,\phi)
&=-(\boldsymbol f,\operatorname{curl}\phi)
&&\forall\phi\in\widetilde\Psi,\\
a_h(\boldsymbol\sigma,\boldsymbol\tau)
+b_h(\boldsymbol\tau,\psi)&=0
&&\forall\boldsymbol\tau\in\widetilde\Sigma.
\end{aligned}
\right.
\]
置
\[
Y=L^r(\Omega)^2,
\qquad
X=\{\boldsymbol\sigma=(\sigma_{ij})\in L^2(\Omega)^4;
\ \sigma_{12}=\sigma_{21}\}\times\Phi_s,
\]
则有 $T\in\mathcal L(Y;X)$; 又因 Stokes 问题正则, 还有
\[
T\in\mathcal L(Y;W^{1,r}(\Omega)^4\times W^{3,r}(\Omega)).
\]
由 \eqref{eq:navier-centered-stream-nonlinearity}, 用
\[
G(\lambda,u)
=-\lambda(\operatorname{tr}(\boldsymbol\sigma)
\operatorname{grad}\psi+\boldsymbol f)
\]
引入非线性对流项. 这是从 $\Lambda\times X$ 到 $L^r(\Omega)^2$
的 $C^\infty$ 映射. 采用这些记号, Navier--Stokes 方程具有标准形式
\begin{equation}\label{eq:navier-centered-gradient-tensor-continuous}
u(\lambda)\in X,
\qquad
F(\lambda,u(\lambda))
\equiv u(\lambda)+TG(\lambda,u(\lambda))=0
\qquad\forall\lambda\in\Lambda.
\end{equation}
关于逼近, 取
\[
X_h=\Sigma_h\times\Phi_h\subset X,
\]
其中
\[
\Sigma_h=\{\boldsymbol\tau\in\widetilde\Sigma;
\ \boldsymbol\tau|_\kappa\in P_{l-1}^4
\quad\forall\kappa\in\mathcal T_h\},
\]
\[
\Phi_h=\{\phi\in\Phi;
\ \phi|_\kappa\in P_l
\quad\forall\kappa\in\mathcal T_h\},
\]
其中整数 $l\geq1$. Stokes 算子离散化为
\[
\left\{
\begin{aligned}
T_h\boldsymbol f&=u_h=(\boldsymbol\sigma_h,\psi_h)\in X_h,\\
b_h(\boldsymbol\sigma_h,\phi_h)
&=-(\boldsymbol f,\operatorname{curl}\phi_h),\\
a_h(\boldsymbol\sigma_h,\boldsymbol\tau_h)
+b_h(\boldsymbol\tau_h,\psi_h)&=0
\end{aligned}
\right\}
\quad
\forall v_h=(\boldsymbol\tau_h,\phi_h)\in X_h,
\]
而 Navier--Stokes 方程逼近为
\begin{equation}\label{eq:navier-centered-gradient-tensor-discrete}
u_h(\lambda)\in X_h,
\qquad
F_h(\lambda,u_h(\lambda))
\equiv u_h(\lambda)+T_hG(\lambda,u_h(\lambda))=0
\qquad\forall\lambda\in\Lambda.
\end{equation}
Subsection~\ref{subsec:incompressible-hhj-additional-results} 的结果给出
$T-T_h$ 的估计
\[
\begin{aligned}
\|(T-T_h)\boldsymbol f\|_X
&=\|\boldsymbol\sigma-\boldsymbol\sigma_h\|_{0,\Omega}
+|\psi-\psi_h|_{1,s,\Omega}\\
&\leq Ch^{2(1-1/r)}|\psi|_{3,r,\Omega}.
\end{aligned}
\]

\MathBookSourcePage{350}
因此, 对分别满足 \eqref{eq:navier-centered-gradient-tensor-continuous} 和
\eqref{eq:navier-centered-gradient-tensor-discrete} 的 $u$ 和 $u_h$,
应用 Theorem~\ref{thm:navier-branch-operator-approximation}、
Theorem~\ref{thm:incompressible-hhj-error} 和
Corollary~\ref{cor:incompressible-hhj-w1s-error}, 得到误差界
\[
\|u(\lambda)-u_h(\lambda)\|_X\leq Ch^k,
\]
只要对某个 $k\in[1,l]$ 有 $\psi(\lambda)\in H^{k+2}(\Omega)$.

最后, 若 $l\geq2$ 时 $\psi(\lambda)\in H^{l+1}(\Omega)$,
或 $l=1$ 时 $\psi(\lambda)\in H^3(\Omega)$,
则与 Theorem~\ref{thm:navier-centered-stream-h1-error} 类似的对偶性论证
给出最优估计
\[
|\psi(\lambda)-\psi_h(\lambda)|_{1,\Omega}\leq Ch^l.
\]
应用 Subsection~\ref{subsec:incompressible-hhj-discontinuous-pressure} 的材料,
还可对压力导出同阶误差界. 证明留作 Exercise.
